\documentclass[11pt]{article}
\usepackage[margin=1in]{geometry}
\usepackage{amsmath,amssymb}
\usepackage{graphicx}
\usepackage{booktabs}
\usepackage{microtype}
\usepackage[hidelinks]{hyperref}
\usepackage{caption}
\captionsetup{font=small,labelfont=bf}

\title{\textbf{Overnight, or Overhyped?}\\[2pt]
\large A Reproducible Audit of the Overnight Return Anomaly\\ and the Market-Manipulation Hypothesis}
\author{Guillain d'Erceville\\
\small\href{https://github.com/Guillain-RDCDE/Open-Alpha-Lab}{github.com/Guillain-RDCDE/Open-Alpha-Lab}}
\date{June 2026}

\begin{document}
\maketitle

\begin{abstract}
\noindent Across world equity markets, close-to-open (``overnight'') returns have
historically dwarfed open-to-close (``intraday'') returns. Knuteson~\cite{knuteson2019,knuteson2023}
reads this as the signature of deliberate manipulation by a large quantitative
fund. We take the empirical fact seriously---and confirm it with
autocorrelation-robust inference and at the individual-firm level---then audit
the inference. We show that (i)~the headline magnitudes are inflated by
compounding, data artefacts, and above all a \emph{calendar-time} confound: the
overnight window spans $\sim$28 hours versus the 6.5-hour day, so per hour the
edge is $\sim$1.3$\times$, not $\sim$4$\times$; (ii)~what survives is largely a
gap-risk premium (beta), sits below realistic trading costs, and has decayed from
a 5-year Sharpe near 2 to near 0.5; (iii)~the strategy is uninvestable at scale,
with square-root market-impact capacity in the single-digit millions of dollars;
and (iv)~the alleged manipulation mechanism is economically self-defeating---a
steelman simulation shows a firm large enough to ``move world markets'' would lose
tens of billions per year running it. A transparent Bayesian analysis returns a
low posterior probability of manipulation that is robust across priors. Every
number is reproducible from the accompanying open-source toolkit.
\end{abstract}

\section{Introduction}
The observation is real and old: U.S.\ equity index returns accrue
overwhelmingly while the market is closed, with the trading day flat to negative
\cite{cooper2008,berkman2012,lou2019}. Knuteson extends it across world markets
and three decades and advances a striking interpretation: a single large
quantitative fund expands its book when the market is illiquid (near the open,
moving prices up) and contracts it when liquid, harvesting the overnight drift it
helps create \cite{knuteson2019,knuteson2023}.

This paper neither dismisses the fact nor accepts the interpretation. We grant
the empirical pattern its fairest hearing, then ask what it actually implies. Our
contribution is an end-to-end, fully reproducible audit that separates three
levels routinely conflated in the popular framing: the \emph{fact} (real), the
\emph{magnitude} (oversold), and the \emph{attribution} (unproven). All results
come from a public toolkit (\texttt{quantlab/}) with a deterministic test suite.

\section{Data and decomposition}
For each trading day $t$ with open/close auction prices, the close-to-close return
splits \emph{exactly} into overnight and intraday legs:
\begin{equation}
(1+r^{\text{on}}_t)(1+r^{\text{id}}_t)=(1+r^{\text{cc}}_t),\qquad
r^{\text{on}}_t=\tfrac{O_t}{C_{t-1}}-1,\;\; r^{\text{id}}_t=\tfrac{C_t}{O_t}-1.
\end{equation}
There is no free parameter. Data are daily OHLC from Yahoo!\ Finance; the
dividend-adjustment mode is treated as a modelling choice (ex-dividend occurs at
the open), with \texttt{split\_only} the default. The test suite verifies the
identity to machine precision ($\sim 10^{-16}$).

\section{The fact is real}
On SPY the overnight mean is $+3.25$ bps/day with a Newey--West (HAC)
$t\approx 4.9$ \cite{neweywest1987}; the intraday mean is insignificant
($t\approx 0.9$). The Lo~\cite{lo2002} Sharpe-ratio $t$-statistic is $\approx 4.5$.
The effect is not an artefact of assuming i.i.d.\ returns.

It is also not a handful of cherry-picked indices. At the \emph{firm} level
across S\&P~500 members ($N\approx 441$ with clean post-2010 history),
$\sim$69\% of stocks have a higher overnight than intraday Sharpe; the median
overnight Sharpe ($\approx 0.56$) exceeds the median intraday Sharpe
($\approx 0.34$); and an equal-weight cross-section earns $+372\%$ overnight
versus $+233\%$ intraday (HAC $t\approx 3.8$), echoing \cite{lou2019}
(Figure~\ref{fig:xs}). Finally, a White~\cite{white2000} Reality
Check over a universe of major indices rejects the null that the best overnight
Sharpe is a product of multiplicity ($p\approx 0$): on liquid indices the effect
survives a data-snooping correction. We therefore concede the fact without
reservation. The interesting questions are about magnitude, cause, and
exploitability.

\begin{figure}[t]
\centering
\includegraphics[width=.8\linewidth]{figures/cross_section.pdf}
\caption{Firm-level overnight vs intraday annualised Sharpe across S\&P~500
members. The overnight distribution is shifted right of the intraday one; dashed
lines mark medians. The effect is a broad cross-sectional fact, not a few-index
fluke. \emph{Caveat:} current membership, hence survivorship bias.}
\label{fig:xs}
\end{figure}

\section{The magnitude is oversold}
Three multipliers inflate the headline.

\paragraph{Compounding on a log axis.} A \emph{constant, innocent} 1 bps/night
drift compounds to $+124\%$ over 32 years; 30 bps/night to trillions of percent.
The explosion is the exponent, not a mechanism.

\paragraph{The clock illusion.} The overnight window is not 6.5 hours like the
day: it runs from the 16:00 close to the 09:30 open (17.5h on weeknights) and
across weekends and holidays (65h$+$), averaging $\sim$28 calendar hours. Placing
both legs on a per-hour footing collapses SPY's overnight advantage from
$\sim$4.1$\times$ \emph{per session} to $\sim$1.3$\times$ \emph{per calendar
hour}. Most of the famous gap is a unit error; a modest genuine premium remains.

\paragraph{Data artefacts.} A few mis-adjusted split/dividend prints mechanically
shovel return from the day into the night; an automatic detector flags them.
These drive the most extreme emerging-market figures, and selecting ``the 25 most
problematic'' from a large universe guarantees significance by construction.

\section{What survives is beta, untradable, and decaying}
Holding the market every night is bearing gap risk continuously. Regressing the
overnight leg on the market (close-to-close) gives $\beta\approx 0.33$: roughly
40\% of SPY's overnight return is a gap-risk premium, not distinct alpha. The
residual alpha ($\approx 1.9$ bps/day) sits \emph{below} the break-even
round-trip cost ($\approx 3.3$ bps): even the non-beta part does not survive
execution. A cost sweep turns a positive gross Sharpe negative at a realistic
5~bps round trip---precisely the fate of the NSPY/NIWM ``night effect'' ETFs
(2022--2023). And the edge is fading: the trailing 5-year overnight Sharpe has
fallen from $\sim$2 (1998) to $\sim$0.5 (2026), the textbook post-publication
decay of \cite{mclean2016} (Figure~\ref{fig:decay}).

\begin{figure}[t]
\centering
\includegraphics[width=.8\linewidth]{figures/alpha_decay.pdf}
\caption{Trailing 5-year annualised Sharpe of the SPY overnight leg. An effect
that fades as it becomes documented is the opposite of an ongoing, deliberate
scheme.}
\label{fig:decay}
\end{figure}

\section{Capacity and the steelman manipulator}
Even a real, free edge is worthless if it cannot hold capital. Under the
square-root market-impact law \cite{almgren2005}, round-trip impact equals SPY's
gross overnight edge at only $\sim\$10$M of capital. We then \emph{steelman} the
manipulation thesis: grant a firm the \emph{entire} overnight drift on its
capital and charge it only round-trip temporary impact (permanent impact nets to
zero over a symmetric round trip, the no-dynamic-arbitrage condition of
\cite{huberman2004,gatheral2010}). Net profit is
\begin{equation}
\Pi(K)=K\Big(r^{\text{on}} - 2\,c\,\sigma\sqrt{K/\text{ADV}\$}\Big),
\end{equation}
positive only for small $K$. The break-even is $\sim\$10$M; at \$1B the firm
loses $\sim\$0.75$B/year, at \$10B $\sim\$25$B/year (Figure~\ref{fig:manip}). A
firm large enough to move world markets could not profit from an edge this thin:
the mechanism is self-defeating at the scale the accusation requires.

\begin{figure}[t]
\centering
\includegraphics[width=.8\linewidth]{figures/manipulator_pnl.pdf}
\caption{Net edge (bps on capital) of the alleged round-trip overnight harvester
versus deployed capital. Profitable only at trivial scale; deeply negative at the
sizes a ``world-moving'' actor would require.}
\label{fig:manip}
\end{figure}

\section{Is it manipulation? A Bayesian reading}
Write the posterior odds for the manipulation hypothesis $M$ against
``risk premium $+$ microstructure'' $S$ given evidence $E_i$:
\begin{equation}
\frac{P(M\mid E)}{P(S\mid E)}=\frac{P(M)}{P(S)}\prod_i \frac{P(E_i\mid M)}{P(E_i\mid S)}.
\end{equation}
The basic night~$>$~day pattern is near-equally likely under both hypotheses, so
it barely discriminates ($\text{LR}\approx 1$). The \emph{discriminating}
evidence---foreign-ETF sign inversions driven by the listing clock, the Chinese
T+1 inversion \cite{qiao2020}, the clock illusion, the negative net P\&L at scale,
and the post-publication decay---all favour $S$. With deliberately moderate,
stated likelihood ratios and a prior of $0.5$, the posterior probability of
manipulation is $\approx 2.4\%$, and it remains low across a wide range of priors
and evidence strengths (Figure~\ref{fig:post}); only a strong prior combined with
deliberately weakened evidence keeps it elevated. The likelihood ratios are
subjective and combined under a generous independence assumption---stated openly,
not hidden---so the robust conclusion is directional: the evidence moves a
reasonable observer \emph{toward} microstructure, not toward fraud.

\begin{figure}[t]
\centering
\includegraphics[width=.8\linewidth]{figures/posterior.pdf}
\caption{Posterior probability of manipulation versus prior, for three evidence
strengths (LR raised to powers $0.5/1/2$). The diagonal is ``posterior = prior''.
Discriminating evidence pulls every curve below the diagonal.}
\label{fig:post}
\end{figure}

\section{Conclusion}
The overnight anomaly is real, broad, and genuinely interesting---and largely
\emph{explained} by ordinary risk premia and market microstructure. Its magnitude
is inflated by compounding and a calendar-time confound; what remains is mostly
beta, is below trading costs, has decayed, and cannot scale. The manipulation
attribution is not supported: the pattern is weak evidence, the discriminating
evidence points the other way, and the alleged mechanism is economically
self-defeating. A beautiful feature of microstructure was dressed up as both a
free lunch and a global fraud; the honest, reproducible answer is more
interesting than either.

\paragraph{Reproducibility.} All figures, tables and numbers are produced by the
open-source toolkit at the URL on the title page: a deterministic, tested
Python package with continuous integration, executed notebooks, and this paper's
figure script (\texttt{paper/make\_figures.py}). \emph{Not investment advice.}

\begin{thebibliography}{99}
\small
\bibitem{almgren2005} Almgren, R., C.\ Thum, E.\ Hauptmann, and H.\ Li. 2005. Direct Estimation of Equity Market Impact. \emph{Risk} 18(7).
\bibitem{berkman2012} Berkman, H., P.\ D.\ Koch, L.\ Tuttle, and Y.\ J.\ Zhang. 2012. Paying Attention: Overnight Returns and the Hidden Cost of Buying at the Open. \emph{JFQA} 47(4): 715--741.
\bibitem{cooper2008} Cooper, M.\ J., M.\ T.\ Cliff, and H.\ Gulen. 2008. Return Differences between Trading and Non-Trading Hours: Like Night and Day. Working paper.
\bibitem{gatheral2010} Gatheral, J. 2010. No-Dynamic-Arbitrage and Market Impact. \emph{Quantitative Finance} 10(7): 749--759.
\bibitem{huberman2004} Huberman, G., and W.\ Stanzl. 2004. Price Manipulation and Quasi-Arbitrage. \emph{Econometrica} 72(4): 1247--1275.
\bibitem{knuteson2019} Knuteson, B. 2019. Celebrating Three Decades of Worldwide Stock Market Manipulation. arXiv:1912.01708.
\bibitem{knuteson2023} Knuteson, B. 2023. Nothing to See Here: How to Say It When You Need to. SSRN 4619084.
\bibitem{lo2002} Lo, A.\ W. 2002. The Statistics of Sharpe Ratios. \emph{Financial Analysts Journal} 58(4): 36--52.
\bibitem{lou2019} Lou, D., C.\ Polk, and S.\ Skouras. 2019. A Tug of War: Overnight versus Intraday Expected Returns. \emph{JFE} 134(1): 192--213.
\bibitem{mclean2016} McLean, R.\ D., and J.\ Pontiff. 2016. Does Academic Research Destroy Stock Return Predictability? \emph{Journal of Finance} 71(1): 5--32.
\bibitem{neweywest1987} Newey, W.\ K., and K.\ D.\ West. 1987. A Simple, Positive Semi-Definite, HAC Covariance Matrix. \emph{Econometrica} 55(3): 703--708.
\bibitem{qiao2020} Qiao, K., and L.\ Dam. 2020. The Overnight Return Puzzle and the ``T+1'' Trading Rule in Chinese Stock Markets. \emph{Journal of Financial Markets} 50: 100534.
\bibitem{white2000} White, H. 2000. A Reality Check for Data Snooping. \emph{Econometrica} 68(5): 1097--1126.
\end{thebibliography}

\end{document}
